% =====================================================================
\section{Organizaci\'on}\label{sec:organizacion}
% =====================================================================

\subsection{Descripci\'on general}
El taller JC\textbar{}cell, ubicado en el cant\'on Quevedo, provincia de Los R\'ios, ofrece
mantenimiento y reparaci\'on de computadoras, laptops y otros dispositivos electr\'onicos.
Entre los servicios que mencion\'o el responsable est\'an el formateo, la instalaci\'on de
sistemas operativos y programas, la limpieza interna, el cambio de componentes y la soluci\'on
de fallas de software y hardware (SE-01, A2; EV-04; Anexo~\ref{anx:registro-entrevista}).

En la entrevista, el responsable estim\'o que el taller lleva cerca de tres a\~nos funcionando
(A2). Tambi\'en explic\'o que trabaja solo: recibe a los clientes, registra los equipos, realiza
los diagn\'osticos y reparaciones, comunica precios y coordina los cobros y las entregas (A1--A4,
B2--B8). Para registrar la recepci\'on anota el tipo, la marca, el modelo, la falla y los datos
de contacto; casi siempre lo hace en una libreta y a veces completa la informaci\'on en su
celular (B2). La antig\"uedad es una estimaci\'on del participante, no una fecha verificada.

\pendiente{Volumen aproximado de equipos atendidos por semana; no se pregunt\'o en SE-01.}

\subsection{Problema identificado}
En SE-01, el responsable describi\'o que consulta sus anotaciones para saber qu\'e equipos est\'an
en revisi\'on, pendientes de autorizaci\'on, en reparaci\'on o terminados; no cuenta con un sistema
que actualice esos estados (B2, B4; EV-04; Anexo~\ref{anx:registro-entrevista}). Cuando se
acumulan varios equipos, buscar entre las notas puede dificultar la identificaci\'on del equipo y
su estado. La principal dificultad que el participante prioriz\'o fue llevar el control de los
equipos que ingresan (C2, C4).

Como caso concreto, relat\'o que un equipo qued\'o esperando un repuesto entre varios trabajos
pendientes. D\'ias despu\'es tuvo que revisar las notas para recordar su estado y qu\'e faltaba hacer;
el tiempo de b\'usqueda le quit\'o tiempo a otras reparaciones. Indic\'o que esto puede ocurrir
cuando se acumulan equipos o una reparaci\'on tarda m\'as de lo previsto, pero no dio una frecuencia
ni un registro con el que medirla (C3). Tambi\'en mencion\'o dificultades para localizar equipos o
piezas, conseguir repuestos, avisar a clientes y recordar cobros pendientes (B2, B5--B7, C2).
Sobre los equipos que a veces no logra ubicar, la entrevista no aclara si se trata de un extrav\'io
temporal o una p\'erdida definitiva. Tampoco se revisaron registros que permitan confirmar esos
casos o establecer cu\'antos ocurren.

Lo anterior corresponde a lo declarado por el responsable en la entrevista. La transcripci\'on se
prepar\'o a partir de las notas del equipo y fue cotejada con la grabaci\'on autorizada (EV-04,
EV-05); aun as\'i, el audio registra el testimonio y no reemplaza los apuntes de recepci\'on o cobro
que permitir\'ian confirmar los casos y medir su frecuencia. El participante propuso revisar
ejemplos anonimizados de recepciones, trabajos y entregas para aclarar la informaci\'on pendiente
(E3). Esos ejemplos a\'un no constan entre las evidencias publicadas, por lo que las cifras de
frecuencia, tiempos de b\'usqueda y pagos olvidados siguen pendientes de contraste (C3, D1, E3).

\subsection{Primer contacto y evidencias}
La Tabla~\ref{tab:evidencias-org} resume las evidencias que respaldan la participaci\'on de la
organizaci\'on. Los originales se custodian en la carpeta institucional privada del equipo y su
registro (nombre de archivo, fecha y huella SHA-256) consta en \texttt{evidencias/MANIFIESTO.md}.

\begin{table}[H]
  \centering
  \caption{Evidencias de la relaci\'on con la organizaci\'on.}
  \label{tab:evidencias-org}
  \noindent\begin{tabularx}{\linewidth}{@{}L{1.2cm}L{3.5cm}L{2.2cm}Y@{}}
    \toprule
    \textbf{ID} & \textbf{Evidencia} & \textbf{Fecha} & \textbf{Descripci\'on} \\
    \midrule
    EV-02 & Registro del primer contacto (conversaci\'on previa)  & \pendiente{fecha} & \pendiente{medio y si existe registro escrito} \\
    EV-03 & Consentimiento informado firmado, PFC-ISR401-CONTX-A03 v1.3 (JC-731) & 20/11/2026 & Responsable del taller; modelo en blanco en el Anexo~\ref{anx:consentimiento} \\
    EV-04 & Transcripci\'on de la entrevista SE-01 (Word) & 21/11/2026 & Notas de la sesi\'on; registro anonimizado en el Anexo~\ref{anx:registro-entrevista} \\
    EV-05 & Grabaci\'on de audio de SE-01 & 21/11/2026 & 12 min 44 s; autorizada en la opci\'on b del consentimiento \\
    \bottomrule
  \end{tabularx}
\end{table}
